\documentclass[fleqn]{article}
\usepackage{amsmath}
\DeclareMathOperator{\sen}{sen} 
\usepackage{amsfonts}
\usepackage{enumitem}
\title{Solución Problemas Fourier Continuo}
\author{Juan Rodríguez}

\setlength{\parindent}{0pt}

\begin{document}
\maketitle
\section{Ejercicio 1}
	Determinar la parte par y la parte impar de la señal escalón unitario

	\[
	u[n]
	\]

	\subsection{Solución}

	La parte par viene dada por:

	\[
	u_{\text{par}}[n]
	=
	\frac{u[n]+u[-n]}{2}
	\]

	Como:

	\[
	u[n]=
	\begin{cases}
	1, & n\geq0 \\[1mm]
	0, & n<0
	\end{cases}
	\]

	se tiene:

	\[
	u[-n]=
	\begin{cases}
	1, & n\leq0 \\[1mm]
	0, & n>0
	\end{cases}
	\]

	Por tanto:

	\[
	\boxed{
	u_{\text{par}}[n]
	=
	\begin{cases}
	1, & n=0 \\[1mm]
	\frac{1}{2}, & n\neq0
	\end{cases}
	}
	\]

	La parte impar viene dada por:

	\[
	u_{\text{impar}}[n]
	=
	\frac{u[n]-u[-n]}{2}
	\]

	Por tanto:

	\[
	\boxed{
	u_{\text{impar}}[n]
	=
	\begin{cases}
	-\frac{1}{2}, & n<0 \\[1mm]
	0, & n=0 \\[1mm]
	\frac{1}{2}, & n>0
	\end{cases}
	}
	\]

\section{Ejercicio 2}
	Obtener, en los casos en que sea posible, el período de las siguientes señales.

	\subsection{Solución}

	Para una señal discreta sinusoidal:

	\[
	x[n]=\cos(\Omega_0 n)
	\quad\text{o}\quad
	x[n]=\sen(\Omega_0 n)
	\]

	es periódica si existe un entero $N>0$ tal que:

	\[
	\Omega_0 N=2\pi k,
	\qquad k\in\mathbb{Z}
	\]

	\textbf{a)}

	\[
	x_1[n]
	=
	\cos\left(\frac{2\pi}{12}n\right)
	\]

	\[
	\frac{2\pi}{12}N
	=
	2\pi k
	\]

	\[
	N=12k
	\]

	El menor valor positivo se obtiene para $k=1$:

	\[
	\boxed{
	N_0=12
	}
	\]

	\textbf{b)}

	\[
	x_2[n]
	=
	\sen\left(\frac{8\pi}{31}n\right)
	\]

	\[
	\frac{8\pi}{31}N
	=
	2\pi k
	\]

	\[
	4N=31k
	\]

	El menor par de enteros que satisface la igualdad es:

	\[
	N=31,
	\qquad
	k=4
	\]

	Por tanto:

	\[
	\boxed{
	N_0=31
	}
	\]

	\textbf{c)}

	\[
	x_3[n]
	=
	\sen\left(\frac{2}{5}n\right)
	\]

	Para que sea periódica debería cumplirse:

	\[
	\frac{2}{5}N
	=
	2\pi k
	\]

	\[
	N=5\pi k
	\]

	Como $\pi$ es irracional, no existen enteros $N>0$ y $k$ que satisfagan la igualdad.

	Por tanto:

	\[
	\boxed{
	x_3[n]\text{ no es periódica}
	}
	\]

	\textbf{d)}

	\[
	x_4[n]
	=
	e^{j\frac{2\pi}{3}n}
	+
	e^{j\frac{3\pi}{4}n}
	\]

	Para el primer término:

	\[
	\frac{2\pi}{3}N_1
	=
	2\pi k
	\]

	\[
	N_1=3
	\]

	Para el segundo término:

	\[
	\frac{3\pi}{4}N_2
	=
	2\pi k
	\]

	\[
	3N_2=8k
	\]

	El menor valor posible es:

	\[
	N_2=8
	\]

	El período fundamental de la suma es el mínimo común múltiplo:

	\[
	N_0
	=
	\operatorname{mcm}(3,8)
	=
	24
	\]

	\[
	\boxed{
	N_0=24
	}
	\]

\section{Ejercicio 3}
	Dada la señal

	\[
	x[n]=
	\begin{cases}
	\left(\frac{1}{2}\right)^n, & n\geq0 \\[1mm]
	0, & n<0
	\end{cases}
	\]

	determinar la energía y la potencia media de $x[n]$.

	\subsection{Solución}

	La energía de una señal discreta es:

	\[
	E
	=
	\sum_{n=-\infty}^{\infty}
	|x[n]|^2
	\]

	Como $x[n]=0$ para $n<0$:

	\[
	E
	=
	\sum_{n=0}^{\infty}
	\left(\frac{1}{2}\right)^{2n}
	\]

	\[
	E
	=
	\sum_{n=0}^{\infty}
	\left(\frac{1}{4}\right)^n
	\]

	Es una serie geométrica de razón:

	\[
	r=\frac{1}{4}
	\]

	Por tanto:

	\[
	E
	=
	\frac{1}{1-\frac{1}{4}}
	=
	\frac{4}{3}
	\]

	\[
	\boxed{
	E=\frac{4}{3}
	}
	\]

	La potencia media es:

	\[
	P
	=
	\lim_{N\to\infty}
	\frac{1}{2N+1}
	\sum_{n=-N}^{N}
	|x[n]|^2
	\]

	\[
	P
	=
	\lim_{N\to\infty}
	\frac{1}{2N+1}
	\sum_{n=0}^{N}
	\left(\frac{1}{4}\right)^n
	\]

	\[
	P
	=
	\lim_{N\to\infty}
	\frac{1}{2N+1}
	\frac{1-\left(\frac{1}{4}\right)^{N+1}}
	{1-\frac{1}{4}}
	\]

	\[
	P=0
	\]

	Por tanto:

	\[
	\boxed{
	P=0
	}
	\]

	\[
	\boxed{
	E=\frac{4}{3},
	\qquad
	P=0
	}
	\]

\section{Ejercicio 4}
	Dada la señal $x[n]$ de la figura, calcular las siguientes señales.

	\subsection{Solución}

	De la gráfica:

	\[
	x[n]=
	\begin{cases}
	\frac{3}{2}, & n=-1 \\[1mm]
	1, & n=0,1,2 \\[1mm]
	\frac{1}{2}, & n=3,4 \\[1mm]
	0, & \text{resto}
	\end{cases}
	\]

	\textbf{a)}

	\[
	y_1[n]=x[n-2]
	\]

	Es un desplazamiento de $2$ posiciones hacia la derecha:

	\[
	\boxed{
	y_1[n]=
	\begin{cases}
	\frac{3}{2}, & n=1 \\[1mm]
	1, & n=2,3,4 \\[1mm]
	\frac{1}{2}, & n=5,6 \\[1mm]
	0, & \text{resto}
	\end{cases}
	}
	\]

	\textbf{b)}

	\[
	y_2[n]=x[4-n]
	\]

	Los valores no nulos son:

	\[
	\boxed{
	y_2[n]=
	\begin{cases}
	\frac{1}{2}, & n=0,1 \\[1mm]
	1, & n=2,3,4 \\[1mm]
	\frac{3}{2}, & n=5 \\[1mm]
	0, & \text{resto}
	\end{cases}
	}
	\]

	\textbf{c)}

	\[
	y_2[n]=x[-n-1]
	\]

	Los valores no nulos son:

	\[
	\boxed{
	y_2[n]=
	\begin{cases}
	\frac{1}{2}, & n=-5,-4 \\[1mm]
	1, & n=-3,-2,-1 \\[1mm]
	\frac{3}{2}, & n=0 \\[1mm]
	0, & \text{resto}
	\end{cases}
	}
	\]

	\textbf{d)}

	\[
	y_3[n]=x[2n]
	\]

	Solo se conservan los valores de $x[n]$ cuyos índices son pares:

	\[
	y_3[0]=x[0]=1
	\]

	\[
	y_3[1]=x[2]=1
	\]

	\[
	y_3[2]=x[4]=\frac{1}{2}
	\]

	Por tanto:

	\[
	\boxed{
	y_3[n]=
	\begin{cases}
	1, & n=0,1 \\[1mm]
	\frac{1}{2}, & n=2 \\[1mm]
	0, & \text{resto}
	\end{cases}
	}
	\]

	\textbf{e)}

	\[
	y_4[n]=x[n]u[2-n]
	\]

	Como:

	\[
	u[2-n]=
	\begin{cases}
	1, & n\leq2 \\[1mm]
	0, & n>2
	\end{cases}
	\]

	se eliminan los valores de $x[n]$ posteriores a $n=2$:

	\[
	\boxed{
	y_4[n]=
	\begin{cases}
	\frac{3}{2}, & n=-1 \\[1mm]
	1, & n=0,1,2 \\[1mm]
	0, & \text{resto}
	\end{cases}
	}
	\]

	\textbf{f)}

	\[
	y_5[n]
	=
	x[n-1]\cdot\delta[n-3]
	\]

	La delta solo es distinta de cero para $n=3$:

	\[
	x[3-1]
	=
	x[2]
	=
	1
	\]

	Por tanto:

	\[
	\boxed{
	y_5[n]=\delta[n-3]
	}
	\]

	\textbf{g)}

	\[
	y_6[n]
	=
	x[n-1]*\delta[n-3]
	\]

	Utilizando la propiedad:

	\[
	f[n]*\delta[n-n_0]
	=
	f[n-n_0]
	\]

	se obtiene:

	\[
	y_6[n]
	=
	x[(n-3)-1]
	=
	x[n-4]
	\]

	Por tanto:

	\[
	\boxed{
	y_6[n]=
	\begin{cases}
	\frac{3}{2}, & n=3 \\[1mm]
	1, & n=4,5,6 \\[1mm]
	\frac{1}{2}, & n=7,8 \\[1mm]
	0, & \text{resto}
	\end{cases}
	}
	\]

\section{Ejercicio 5}
	Obtener los coeficientes de la serie de Fourier de la señal discreta periódica $x[n]$ de la figura.

	\subsection{Solución}

	De la gráfica se observa que:

	\[
	N_0=4
	\]

	y un período de la señal es:

	\[
	x[0]=4,
	\qquad
	x[1]=1,
	\qquad
	x[2]=2,
	\qquad
	x[3]=3
	\]

	La frecuencia fundamental es:

	\[
	\Omega_0
	=
	\frac{2\pi}{N_0}
	=
	\frac{\pi}{2}
	\]

	Los coeficientes vienen dados por:

	\[
	c_k
	=
	\frac{1}{N_0}
	\sum_{n=0}^{N_0-1}
	x[n]e^{-jk\Omega_0n}
	\]

	Para $k=0$:

	\[
	c_0
	=
	\frac{1}{4}
	(4+1+2+3)
	=
	\frac{5}{2}
	\]

	Para $k=1$:

	\[
	c_1
	=
	\frac{1}{4}
	\left(
	4
	+
	e^{-j\pi/2}
	+
	2e^{-j\pi}
	+
	3e^{-j3\pi/2}
	\right)
	\]

	\[
	c_1
	=
	\frac{1}{4}
	\left(
	4-j-2+3j
	\right)
	\]

	\[
	c_1
	=
	\frac{1+j}{2}
	\]

	Para $k=2$:

	\[
	c_2
	=
	\frac{1}{4}
	\left(
	4
	+
	e^{-j\pi}
	+
	2e^{-j2\pi}
	+
	3e^{-j3\pi}
	\right)
	\]

	\[
	c_2
	=
	\frac{1}{4}
	(4-1+2-3)
	=
	\frac{1}{2}
	\]

	Para $k=3$:

	\[
	c_3
	=
	\frac{1}{4}
	\left(
	4
	+
	e^{-j3\pi/2}
	+
	2e^{-j3\pi}
	+
	3e^{-j9\pi/2}
	\right)
	\]

	\[
	c_3
	=
	\frac{1}{4}
	\left(
	4+j-2-3j
	\right)
	\]

	\[
	c_3
	=
	\frac{1-j}{2}
	\]

	Por tanto:

	\[
	\boxed{
	c_0=\frac{5}{2},
	\qquad
	c_1=\frac{1+j}{2},
	\qquad
	c_2=\frac{1}{2},
	\qquad
	c_3=\frac{1-j}{2}
	}
	\]

	Los coeficientes son periódicos de período $4$:

	\[
	\boxed{
	c_{k+4}=c_k
	}
	\]

\section{Ejercicio 6}
	Desarrollar en serie de Fourier la señal discreta periódica $x[n]$ de la figura.

	\subsection{Solución}

	De la gráfica se observa que:

	\[
	N_0=6
	\]

	y un período de la señal es:

	\[
	x[0]=1,\qquad
	x[1]=1,\qquad
	x[2]=0,\qquad
	x[3]=0,\qquad
	x[4]=0,\qquad
	x[5]=1
	\]

	La frecuencia fundamental es:

	\[
	\Omega_0
	=
	\frac{2\pi}{N_0}
	=
	\frac{\pi}{3}
	\]

	Los coeficientes vienen dados por:

	\[
	c_k
	=
	\frac{1}{6}
	\sum_{n=0}^{5}
	x[n]e^{-jk\Omega_0n}
	\]

	Como solo son distintos de cero los términos $n=0,1,5$:

	\[
	c_k
	=
	\frac{1}{6}
	\left(
	1+
	e^{-jk\pi/3}
	+
	e^{-j5k\pi/3}
	\right)
	\]

	Como:

	\[
	e^{-j5k\pi/3}
	=
	e^{jk\pi/3}
	\]

	se tiene:

	\[
	c_k
	=
	\frac{1}{6}
	\left(
	1+
	e^{-jk\pi/3}
	+
	e^{jk\pi/3}
	\right)
	\]

	\[
	\boxed{
	c_k
	=
	\frac{1+2\cos\left(\frac{k\pi}{3}\right)}{6}
	}
	\]

	Calculamos los coeficientes de un período:

	\[
	c_0=\frac{1}{2}
	\]

	\[
	c_1=\frac{1}{3}
	\]

	\[
	c_2=0
	\]

	\[
	c_3=-\frac{1}{6}
	\]

	\[
	c_4=0
	\]

	\[
	c_5=\frac{1}{3}
	\]

	Por tanto:

	\[
	\boxed{
	c_0=\frac{1}{2},
	\qquad
	c_1=\frac{1}{3},
	\qquad
	c_2=0,
	\qquad
	c_3=-\frac{1}{6},
	\qquad
	c_4=0,
	\qquad
	c_5=\frac{1}{3}
	}
	\]

	y el desarrollo en serie de Fourier es:

	\[
	\boxed{
	x[n]
	=
	\sum_{k=0}^{5}
	c_k e^{jk\frac{\pi}{3}n}
	}
	\]

\section{Ejercicio 7}
	Obtener el valor medio y la potencia media de las señales periódicas de los problemas 5 y 6.

	\subsection{Solución}

	Para la señal del ejercicio 5:

	\[
	x[n]=\{4,1,2,3\},
	\qquad
	N_0=4
	\]

	El valor medio es:

	\[
	\overline{x}
	=
	\frac{1}{4}(4+1+2+3)
	\]

	\[
	\boxed{
	\overline{x}=\frac{5}{2}
	}
	\]

	La potencia media es:

	\[
	P
	=
	\frac{1}{4}
	\left(
	4^2+1^2+2^2+3^2
	\right)
	\]

	\[
	P
	=
	\frac{30}{4}
	\]

	\[
	\boxed{
	P=\frac{15}{2}
	}
	\]

	Para la señal del ejercicio 6:

	\[
	x[n]=\{1,1,0,0,0,1\},
	\qquad
	N_0=6
	\]

	El valor medio es:

	\[
	\overline{x}
	=
	\frac{1}{6}(1+1+1)
	\]

	\[
	\boxed{
	\overline{x}=\frac{1}{2}
	}
	\]

	La potencia media es:

	\[
	P
	=
	\frac{1}{6}
	\left(
	1^2+1^2+1^2
	\right)
	\]

	\[
	\boxed{
	P=\frac{1}{2}
	}
	\]

\section{Ejercicio 8}
	Determinar los coeficientes del desarrollo en serie de Fourier de las siguientes señales.

	\subsection{Solución}

	\textbf{a)}

	\[
	x_1[n]
	=
	\cos\left(\frac{\pi}{4}n\right)
	\]

	Como:

	\[
	\frac{\pi}{4}
	=
	\frac{2\pi}{8}
	\]

	se tiene:

	\[
	N_0=8,
	\qquad
	\Omega_0=\frac{\pi}{4}
	\]

	Utilizando:

	\[
	\cos(\Omega_0n)
	=
	\frac{1}{2}e^{j\Omega_0n}
	+
	\frac{1}{2}e^{-j\Omega_0n}
	\]

	se obtiene:

	\[
	\boxed{
	c_1=\frac{1}{2},
	\qquad
	c_7=\frac{1}{2}
	}
	\]

	y el resto de coeficientes son nulos.

	\textbf{b)}

	\[
	x_2[n]
	=
	\cos\left(\frac{\pi}{3}n\right)
	+
	\sen\left(\frac{\pi}{4}n\right)
	\]

	Los períodos de ambas señales son:

	\[
	N_1=6,
	\qquad
	N_2=8
	\]

	Por tanto:

	\[
	N_0
	=
	\operatorname{mcm}(6,8)
	=
	24
	\]

	\[
	\Omega_0
	=
	\frac{2\pi}{24}
	=
	\frac{\pi}{12}
	\]

	Entonces:

	\[
	\frac{\pi}{3}=4\Omega_0
	\]

	y:

	\[
	\frac{\pi}{4}=3\Omega_0
	\]

	Por tanto:

	\[
	\cos(4\Omega_0n)
	=
	\frac{1}{2}e^{j4\Omega_0n}
	+
	\frac{1}{2}e^{-j4\Omega_0n}
	\]

	y:

	\[
	\sen(3\Omega_0n)
	=
	\frac{1}{2j}e^{j3\Omega_0n}
	-
	\frac{1}{2j}e^{-j3\Omega_0n}
	\]

	Así:

	\[
	\boxed{
	c_4=\frac{1}{2},
	\qquad
	c_{20}=\frac{1}{2},
	\qquad
	c_3=-\frac{j}{2},
	\qquad
	c_{21}=\frac{j}{2}
	}
	\]

	y el resto son nulos.

	\textbf{c)}

	\[
	x_3[n]
	=
	\cos^2\left(\frac{\pi}{8}n\right)
	\]

	Utilizamos:

	\[
	\cos^2(\alpha)
	=
	\frac{1+\cos(2\alpha)}{2}
	\]

	\[
	x_3[n]
	=
	\frac{1}{2}
	+
	\frac{1}{2}
	\cos\left(\frac{\pi}{4}n\right)
	\]

	El período fundamental es:

	\[
	N_0=8
	\]

	y:

	\[
	\frac{1}{2}
	\cos\left(\frac{\pi}{4}n\right)
	=
	\frac{1}{4}e^{j\pi n/4}
	+
	\frac{1}{4}e^{-j\pi n/4}
	\]

	Por tanto:

	\[
	\boxed{
	c_0=\frac{1}{2},
	\qquad
	c_1=\frac{1}{4},
	\qquad
	c_7=\frac{1}{4}
	}
	\]

	y el resto son nulos.

\section{Ejercicio 9}
	Una señal discreta periódica $x[n]$ de período $N_0=4$ tiene los coeficientes:

	\[
	\{3,2-j,4+3j,-9\}
	\]

	Determinar los coeficientes de:

	\[
	y[n]
	=
	x[n]e^{j\frac{3\pi}{2}n}
	\]

	\subsection{Solución}

	Como:

	\[
	N_0=4
	\]

	se tiene:

	\[
	\Omega_0
	=
	\frac{2\pi}{4}
	=
	\frac{\pi}{2}
	\]

	Por tanto:

	\[
	e^{j\frac{3\pi}{2}n}
	=
	e^{j3\Omega_0n}
	\]

	Por la propiedad de desplazamiento en frecuencia:

	\[
	y[n]
	=
	x[n]e^{j3\Omega_0n}
	\]

	implica:

	\[
	d_k=c_{k-3}
	\]

	Como los coeficientes son periódicos de período $4$:

	\[
	d_0=c_1=2-j
	\]

	\[
	d_1=c_2=4+3j
	\]

	\[
	d_2=c_3=-9
	\]

	\[
	d_3=c_0=3
	\]

	Por tanto:

	\[
	\boxed{
	\{d_k\}
	=
	\{2-j,4+3j,-9,3\}
	}
	\]

\section{Ejercicio 10}
	Calcular la convolución

	\[
	z[n]=x[n]*y[n]
	\]

	dadas las señales:

	\[
	x[n]
	=
	2\delta[n]
	+
	\delta[n-2]
	-
	\delta[n-3]
	\]

	\[
	y[n]
	=
	\delta[n+1]
	+
	2\delta[n]
	-
	3\delta[n-1]
	\]

	\subsection{Solución}

	Utilizamos:

	\[
	\delta[n-a]*\delta[n-b]
	=
	\delta[n-(a+b)]
	\]

	Distribuyendo la convolución:

	\[
	z[n]
	=
	2\delta[n]*
	\left(
	\delta[n+1]+2\delta[n]-3\delta[n-1]
	\right)
	\]

	\[
	+
	\delta[n-2]*
	\left(
	\delta[n+1]+2\delta[n]-3\delta[n-1]
	\right)
	\]

	\[
	-
	\delta[n-3]*
	\left(
	\delta[n+1]+2\delta[n]-3\delta[n-1]
	\right)
	\]

	Se obtiene:

	\[
	z[n]
	=
	2\delta[n+1]
	+
	4\delta[n]
	-
	6\delta[n-1]
	\]

	\[
	+
	\delta[n-1]
	+
	2\delta[n-2]
	-
	3\delta[n-3]
	\]

	\[
	-
	\delta[n-2]
	-
	2\delta[n-3]
	+
	3\delta[n-4]
	\]

	Agrupando términos:

	\[
	\boxed{
	z[n]
	=
	2\delta[n+1]
	+
	4\delta[n]
	-
	5\delta[n-1]
	+
	\delta[n-2]
	-
	5\delta[n-3]
	+
	3\delta[n-4]
	}
	\]

\section{Ejercicio 11}
	Calcular la convolución lineal y circular de:

	\[
	x[n]
	=
	\delta[n]
	+
	2\delta[n-1]
	+
	3\delta[n-2]
	+
	3\delta[n-3]
	\]

	\[
	y[n]
	=
	\delta[n]
	-
	\delta[n-1]
	+
	\delta[n-2]
	\]

	\subsection{Solución}

	Las secuencias son:

	\[
	x[n]=\{1,2,3,3\}
	\]

	\[
	y[n]=\{1,-1,1\}
	\]

	Para la convolución lineal:

	\[
	z[n]
	=
	\sum_{k=-\infty}^{\infty}
	x[k]y[n-k]
	\]

	Calculamos los valores:

	\[
	z[0]=1
	\]

	\[
	z[1]=-1+2=1
	\]

	\[
	z[2]=1-2+3=2
	\]

	\[
	z[3]=2-3+3=2
	\]

	\[
	z[4]=3-3=0
	\]

	\[
	z[5]=3
	\]

	Por tanto:

	\[
	\boxed{
	z[n]=\{1,1,2,2,0,3\}
	}
	\]

	Para la convolución circular tomamos:

	\[
	N=\max(4,3)=4
	\]

	y completamos:

	\[
	y[n]=\{1,-1,1,0\}
	\]

	La convolución circular puede obtenerse agrupando los términos de la convolución lineal módulo $4$:

	\[
	z_c[0]
	=
	z[0]+z[4]
	=
	1
	\]

	\[
	z_c[1]
	=
	z[1]+z[5]
	=
	4
	\]

	\[
	z_c[2]
	=
	z[2]
	=
	2
	\]

	\[
	z_c[3]
	=
	z[3]
	=
	2
	\]

	Por tanto:

	\[
	\boxed{
	z_c[n]
	=
	\{1,4,2,2\}
	}
	\]

\section{Ejercicio 12}
	Calcular la convolución circular

	\[
	x[n]\overset{10}{*}y[n]
	\]

	donde:

	\[
	x[n]=
	\begin{cases}
	3n, & 0\leq n\leq6 \\[1mm]
	0, & \text{resto}
	\end{cases}
	\]

	y:

	\[
	y[n]=
	\begin{cases}
	\displaystyle\frac{n}{3}, & 0\leq n\leq5 \\[1mm]
	0, & \text{resto}
	\end{cases}
	\]

	\subsection{Solución}

	Como la convolución circular es de longitud $N=10$, escribimos:

	\[
	x[n]
	=
	\{0,3,6,9,12,15,18,0,0,0\}
	\]

	\[
	y[n]
	=
	\left\{
	0,\frac{1}{3},\frac{2}{3},1,
	\frac{4}{3},\frac{5}{3},0,0,0,0
	\right\}
	\]

	La convolución circular viene dada por:

	\[
	z[n]
	=
	\sum_{k=0}^{9}
	x[k]y[(n-k)\operatorname{mod}10]
	\]

	Calculando para $n=0,\ldots,9$ se obtiene:

	\[
	z[0]=49
	\]

	\[
	z[1]=30
	\]

	\[
	z[2]=1
	\]

	\[
	z[3]=4
	\]

	\[
	z[4]=10
	\]

	\[
	z[5]=20
	\]

	\[
	z[6]=35
	\]

	\[
	z[7]=50
	\]

	\[
	z[8]=58
	\]

	\[
	z[9]=58
	\]

	Por tanto:

	\[
	\boxed{
	z[n]
	=
	\{49,30,1,4,10,20,35,50,58,58\}
	}
	\]

\section{Ejercicio 13}
	Determinar la DTFT de la señal

	\[
	x[n]=a^nu[n],
	\qquad |a|<1
	\]

	\subsection{Solución}

	La DTFT viene dada por:

	\[
	X(\Omega)
	=
	\sum_{n=-\infty}^{\infty}
	x[n]e^{-j\Omega n}
	\]

	Como $u[n]=0$ para $n<0$:

	\[
	X(\Omega)
	=
	\sum_{n=0}^{\infty}
	a^ne^{-j\Omega n}
	\]

	\[
	X(\Omega)
	=
	\sum_{n=0}^{\infty}
	\left(
	ae^{-j\Omega}
	\right)^n
	\]

	Es una serie geométrica de razón:

	\[
	r=ae^{-j\Omega}
	\]

	Como:

	\[
	|r|=|a|<1
	\]

	se tiene:

	\[
	X(\Omega)
	=
	\frac{1}{1-ae^{-j\Omega}}
	\]

	Por tanto:

	\[
	\boxed{
	X(\Omega)
	=
	\frac{1}{1-ae^{-j\Omega}}
	}
	\]

\section{Ejercicio 14}
	Determinar la DTFT de la señal

	\[
	x[n]=a^{|n|},
	\qquad |a|<1
	\]

	\subsection{Solución}

	Aplicamos la definición:

	\[
	X(\Omega)
	=
	\sum_{n=-\infty}^{\infty}
	a^{|n|}e^{-j\Omega n}
	\]

	Separamos:

	\[
	X(\Omega)
	=
	1
	+
	\sum_{n=1}^{\infty}
	a^ne^{-j\Omega n}
	+
	\sum_{n=-\infty}^{-1}
	a^{|n|}e^{-j\Omega n}
	\]

	Haciendo $m=-n$ en la última suma:

	\[
	X(\Omega)
	=
	1
	+
	\sum_{n=1}^{\infty}
	\left(ae^{-j\Omega}\right)^n
	+
	\sum_{n=1}^{\infty}
	\left(ae^{j\Omega}\right)^n
	\]

	Utilizando la suma de la serie geométrica:

	\[
	X(\Omega)
	=
	1
	+
	\frac{ae^{-j\Omega}}
	{1-ae^{-j\Omega}}
	+
	\frac{ae^{j\Omega}}
	{1-ae^{j\Omega}}
	\]

	Operando:

	\[
	X(\Omega)
	=
	\frac{1-a^2}
	{
	(1-ae^{-j\Omega})(1-ae^{j\Omega})
	}
	\]

	Como:

	\[
	e^{j\Omega}+e^{-j\Omega}
	=
	2\cos(\Omega)
	\]

	se obtiene:

	\[
	\boxed{
	X(\Omega)
	=
	\frac{1-a^2}
	{1-2a\cos(\Omega)+a^2}
	}
	\]

\section{Ejercicio 15}
	Determinar la DTFT de las señales

	\[
	x[n]
	=
	a^nu[n]\sen(\Omega_0n)
	\]

	e:

	\[
	y[n]
	=
	a^nu[n]*\sen(\Omega_0n),
	\qquad |a|<1
	\]

	\subsection{Solución}

	Para $x[n]$, utilizamos:

	\[
	\sen(\Omega_0n)
	=
	\frac{
	e^{j\Omega_0n}
	-
	e^{-j\Omega_0n}
	}{2j}
	\]

	Entonces:

	\[
	x[n]
	=
	\frac{1}{2j}
	\left[
	a^nu[n]e^{j\Omega_0n}
	-
	a^nu[n]e^{-j\Omega_0n}
	\right]
	\]

	Sabemos que:

	\[
	a^nu[n]
	\overset{\mathcal{F}}{\longleftrightarrow}
	\frac{1}{1-ae^{-j\Omega}}
	\]

	Aplicando la propiedad de modulación:

	\[
	\boxed{
	X(\Omega)
	=
	\frac{1}{2j}
	\left[
	\frac{1}
	{1-ae^{-j(\Omega-\Omega_0)}}
	-
	\frac{1}
	{1-ae^{-j(\Omega+\Omega_0)}}
	\right]
	}
	\]

	Para $y[n]$, utilizamos la propiedad de la convolución:

	\[
	y[n]
	=
	a^nu[n]*\sen(\Omega_0n)
	\]

	\[
	Y(\Omega)
	=
	A(\Omega)S(\Omega)
	\]

	donde:

	\[
	A(\Omega)
	=
	\frac{1}{1-ae^{-j\Omega}}
	\]

	y:

	\[
	S(\Omega)
	=
	\frac{\pi}{j}
	\left[
	\delta_p(\Omega-\Omega_0)
	-
	\delta_p(\Omega+\Omega_0)
	\right]
	\]

	Por tanto:

	\[
	Y(\Omega)
	=
	\frac{\pi}{j}
	\left[
	\frac{
	\delta_p(\Omega-\Omega_0)
	}{
	1-ae^{-j\Omega_0}
	}
	-
	\frac{
	\delta_p(\Omega+\Omega_0)
	}{
	1-ae^{j\Omega_0}
	}
	\right]
	\]

	\[
	\boxed{
	Y(\Omega)
	=
	\frac{\pi}{j}
	\left[
	\frac{
	\delta_p(\Omega-\Omega_0)
	}{
	1-ae^{-j\Omega_0}
	}
	-
	\frac{
	\delta_p(\Omega+\Omega_0)
	}{
	1-ae^{j\Omega_0}
	}
	\right]
	}
	\]

\section{Ejercicio 16}
	Determinar la DTFT de la señal

	\[
	x[n]=
	\begin{cases}
	1, & |n|\leq N \\[1mm]
	0, & \text{resto}
	\end{cases}
	\]

	\subsection{Solución}

	Aplicamos la definición:

	\[
	X(\Omega)
	=
	\sum_{n=-\infty}^{\infty}
	x[n]e^{-j\Omega n}
	\]

	Como la señal solo es distinta de cero para $-N\leq n\leq N$:

	\[
	X(\Omega)
	=
	\sum_{n=-N}^{N}
	e^{-j\Omega n}
	\]

	Escribimos:

	\[
	X(\Omega)
	=
	e^{jN\Omega}
	\sum_{k=0}^{2N}
	e^{-j\Omega k}
	\]

	Utilizando la suma geométrica:

	\[
	X(\Omega)
	=
	e^{jN\Omega}
	\frac{
	1-e^{-j(2N+1)\Omega}
	}{
	1-e^{-j\Omega}
	}
	\]

	Utilizamos:

	\[
	1-e^{-j\alpha}
	=
	2je^{-j\alpha/2}
	\sen\left(\frac{\alpha}{2}\right)
	\]

	Entonces:

	\[
	X(\Omega)
	=
	e^{jN\Omega}
	\frac{
	2je^{-j(2N+1)\Omega/2}
	\sen\left(\frac{(2N+1)\Omega}{2}\right)
	}{
	2je^{-j\Omega/2}
	\sen\left(\frac{\Omega}{2}\right)
	}
	\]

	Los términos exponenciales se cancelan y se obtiene:

	\[
	\boxed{
	X(\Omega)
	=
	\frac{
	\sen\left(\frac{(2N+1)\Omega}{2}\right)
	}{
	\sen\left(\frac{\Omega}{2}\right)
	}
	}
	\]

	En $\Omega=2\pi k$:

	\[
	\boxed{
	X(2\pi k)=2N+1
	}
	\]

\section{Ejercicio 17}
	Calcular la DTFT inversa de

	\[
	X(\Omega)=
	\begin{cases}
	1, & |\Omega|\leq W<\pi \\[1mm]
	0, & \text{resto}
	\end{cases}
	\]

	y representar gráficamente tanto $X(\Omega)$ como $x[n]$.

	\subsection{Solución}

	La DTFT inversa viene dada por:

	\[
	x[n]
	=
	\frac{1}{2\pi}
	\int_{-\pi}^{\pi}
	X(\Omega)e^{j\Omega n}\,d\Omega
	\]

	Como $X(\Omega)=1$ para $|\Omega|\leq W$:

	\[
	x[n]
	=
	\frac{1}{2\pi}
	\int_{-W}^{W}
	e^{j\Omega n}\,d\Omega
	\]

	Si $n\neq0$:

	\[
	x[n]
	=
	\frac{1}{2\pi}
	\left[
	\frac{e^{j\Omega n}}{jn}
	\right]_{-W}^{W}
	\]

	\[
	x[n]
	=
	\frac{1}{2\pi}
	\frac{
	e^{jWn}-e^{-jWn}
	}{jn}
	\]

	Utilizando:

	\[
	e^{jWn}-e^{-jWn}
	=
	2j\sen(Wn)
	\]

	se obtiene:

	\[
	x[n]
	=
	\frac{\sen(Wn)}{\pi n}
	\]

	Para $n=0$:

	\[
	x[0]
	=
	\frac{1}{2\pi}
	\int_{-W}^{W}d\Omega
	=
	\frac{W}{\pi}
	\]

	Por tanto:

	\[
	\boxed{
	x[n]=
	\begin{cases}
	\displaystyle\frac{W}{\pi}, & n=0 \\[3mm]
	\displaystyle\frac{\sen(Wn)}{\pi n}, & n\neq0
	\end{cases}
	}
	\]

	La señal $X(\Omega)$ es un pulso rectangular entre $-W$ y $W$,
	periódico con período $2\pi$, mientras que $x[n]$ tiene forma de sinc discreta.

\section{Ejercicio 18}
	Calcular la DFT de la secuencia

	\[
	x[n]=[2,0,-1,3]
	\]

	\subsection{Solución}

	La longitud de la secuencia es:

	\[
	N=4
	\]

	La DFT viene dada por:

	\[
	X[k]
	=
	\sum_{n=0}^{3}
	x[n]e^{-j\frac{2\pi}{4}kn}
	\]

	Para $k=0$:

	\[
	X[0]
	=
	2+0-1+3
	=
	4
	\]

	Para $k=1$:

	\[
	X[1]
	=
	2
	-
	e^{-j\pi}
	+
	3e^{-j3\pi/2}
	\]

	\[
	X[1]
	=
	2+1+3j
	\]

	\[
	X[1]=3+3j
	\]

	Para $k=2$:

	\[
	X[2]
	=
	2
	-
	e^{-j2\pi}
	+
	3e^{-j3\pi}
	\]

	\[
	X[2]
	=
	2-1-3
	=
	-2
	\]

	Como la secuencia es real:

	\[
	X[3]=X^*[1]
	\]

	\[
	X[3]=3-3j
	\]

	Por tanto:

	\[
	\boxed{
	X[k]
	=
	[4,\;3+3j,\;-2,\;3-3j]
	}
	\]

\section{Ejercicio 19}
	Calcular la DFT inversa de la secuencia

	\[
	X[k]
	=
	[0,-3-3j,-2,-3+3j]
	\]

	\subsection{Solución}

	La longitud es:

	\[
	N=4
	\]

	La DFT inversa viene dada por:

	\[
	x[n]
	=
	\frac{1}{4}
	\sum_{k=0}^{3}
	X[k]e^{j\frac{2\pi}{4}kn}
	\]

	Para $n=0$:

	\[
	x[0]
	=
	\frac{1}{4}
	\left[
	(-3-3j)-2+(-3+3j)
	\right]
	\]

	\[
	x[0]=-2
	\]

	Para $n=1$:

	\[
	x[1]
	=
	\frac{1}{4}
	\left[
	(-3-3j)j
	+
	(-2)(-1)
	+
	(-3+3j)(-j)
	\right]
	\]

	\[
	x[1]=2
	\]

	Para $n=2$:

	\[
	x[2]
	=
	\frac{1}{4}
	\left[
	(-3-3j)(-1)
	-2
	+
	(-3+3j)(-1)
	\right]
	\]

	\[
	x[2]=1
	\]

	Para $n=3$:

	\[
	x[3]
	=
	\frac{1}{4}
	\left[
	(-3-3j)(-j)
	+
	(-2)(-1)
	+
	(-3+3j)j
	\right]
	\]

	\[
	x[3]=-1
	\]

	Por tanto:

	\[
	\boxed{
	x[n]=[-2,2,1,-1]
	}
	\]

\section{Ejercicio 20}
	Proporcionar un ejemplo de secuencia simétrica y otro de secuencia antisimétrica
	para $N=5$ y para $N=6$.

	\subsection{Solución}

	Una secuencia es simétrica si:

	\[
	x[n]=x[N-n]
	\]

	y es antisimétrica si:

	\[
	x[n]=-x[N-n]
	\]

	Para $N=5$, un ejemplo de secuencia simétrica es:

	\[
	\boxed{
	x[n]=[1,2,3,3,2]
	}
	\]

	ya que:

	\[
	x[1]=x[4],
	\qquad
	x[2]=x[3]
	\]

	Un ejemplo de secuencia antisimétrica es:

	\[
	\boxed{
	x[n]=[0,1,2,-2,-1]
	}
	\]

	Para $N=6$, un ejemplo de secuencia simétrica es:

	\[
	\boxed{
	x[n]=[1,2,3,4,3,2]
	}
	\]

	ya que:

	\[
	x[1]=x[5],
	\qquad
	x[2]=x[4]
	\]

	Un ejemplo de secuencia antisimétrica es:

	\[
	\boxed{
	x[n]=[0,1,2,0,-2,-1]
	}
	\]

	En este caso:

	\[
	x[0]=0,
	\qquad
	x[3]=0
	\]

	ya que estos índices son simétricos consigo mismos.

\section{Ejercicio 21}
	Dada la transformación $X[k]$ asociada a una secuencia real $x[n]$ con $N=14$:

	\[
	X[0]=12
	\]

	\[
	X[1]=-1+3j,\quad
	X[2]=3+4j,\quad
	X[3]=-1+5j
	\]

	\[
	X[4]=-2+2j,\quad
	X[5]=6+3j,\quad
	X[6]=-2-3j,\quad
	X[7]=10
	\]

	calcular el resto de coeficientes.

	\subsection{Solución}

	Como $x[n]$ es una secuencia real, se cumple:

	\[
	X[k]
	=
	X^*[N-k]
	\]

	Por tanto:

	\[
	X[8]
	=
	X^*[6]
	=
-2+3j
	\]

	\[
	X[9]
	=
	X^*[5]
	=
	6-3j
	\]

	\[
	X[10]
	=
	X^*[4]
	=
	-2-2j
	\]

	\[
	X[11]
	=
	X^*[3]
	=
	-1-5j
	\]

	\[
	X[12]
	=
	X^*[2]
	=
	3-4j
	\]

	\[
	X[13]
	=
	X^*[1]
	=
	-1-3j
	\]

	Por tanto:

	\[
	\boxed{
	\begin{aligned}
	X[8]&=-2+3j \\
	X[9]&=6-3j \\
	X[10]&=-2-2j \\
	X[11]&=-1-5j \\
	X[12]&=3-4j \\
	X[13]&=-1-3j
	\end{aligned}
	}
	\]

\section{Ejercicio 22}
	Dada una secuencia de números reales $x[n]$ de longitud $N=9$, se sabe que su transformada DFT es:

	\[
	X[k]
	=
	[3.1,\;a,\;2.5+4.6j,\;b,\;-1.7+5.2j,\;c,\;9.3+6.3j,\;d,\;5.5-8j]
	\]

	Determinar el valor de las constantes $a$, $b$, $c$ y $d$.

	\subsection{Solución}

	Como $x[n]$ es una secuencia real, se cumple:

	\[
	X[k]
	=
	X^*[N-k]
	\]

	Por tanto:

	\[
	X[1]
	=
	X^*[8]
	\]

	\[
	a
	=
	5.5+8j
	\]

	Además:

	\[
	X[2]
	=
	X^*[7]
	\]

	\[
	d
	=
	2.5-4.6j
	\]

	Para $k=3$:

	\[
	X[3]
	=
	X^*[6]
	\]

	\[
	b
	=
	9.3-6.3j
	\]

	Y para $k=4$:

	\[
	X[4]
	=
	X^*[5]
	\]

	\[
	c
	=
	-1.7-5.2j
	\]

	Por tanto:

	\[
	\boxed{
	\begin{aligned}
	a&=5.5+8j\\
	b&=9.3-6.3j\\
	c&=-1.7-5.2j\\
	d&=2.5-4.6j
	\end{aligned}
	}
	\]

\section{Ejercicio 23}
	Realizar de nuevo la convolución circular del problema 11 calculando previamente la DFT de cada secuencia.

	\subsection{Solución}

	Las secuencias del ejercicio 11 son:

	\[
	x[n]=[1,2,3,3]
	\]

	y:

	\[
	y[n]=[1,-1,1,0]
	\]

	con:

	\[
	N=4
	\]

	Calculamos la DFT de $x[n]$:

	\[
	X[k]
	=
	\sum_{n=0}^{3}
	x[n]e^{-j\frac{2\pi}{4}kn}
	\]

	Se obtiene:

	\[
	\boxed{
	X[k]=[9,-2+j,-1,-2-j]
	}
	\]

	Para $y[n]$:

	\[
	Y[k]
	=
	\sum_{n=0}^{3}
	y[n]e^{-j\frac{2\pi}{4}kn}
	\]

	Se obtiene:

	\[
	\boxed{
	Y[k]=[1,j,3,-j]
	}
	\]

	La DFT de la convolución circular es:

	\[
	Z[k]
	=
	X[k]Y[k]
	\]

	Por tanto:

	\[
	Z[k]
	=
	[9,-2+j,-1,-2-j]
	\cdot
	[1,j,3,-j]
	\]

	\[
	\boxed{
	Z[k]
	=
	[9,-1-2j,-3,-1+2j]
	}
	\]

	Aplicando la DFT inversa:

	\[
	z[n]
	=
	\frac{1}{4}
	\sum_{k=0}^{3}
	Z[k]e^{j\frac{2\pi}{4}kn}
	\]

	se obtiene:

	\[
	\boxed{
	z[n]=[1,4,2,2]
	}
	\]

\section{Ejercicio 24}
	Dadas las siguientes secuencias reales, determinar cuáles tendrán una transformada DFT cuyos elementos sean imaginarios puros:

	\[
	\text{a) }x[n]=[1,1,1,0,0,0,1,1]
	\]

	\[
	\text{b) }x[n]=[1,1,0,0,0,0,-1,-1]
	\]

	\[
	\text{c) }x[n]=[0,1,1,0,0,0,-1,-1]
	\]

	\[
	\text{d) }x[n]=[0,1,1,0,0,0,1,1]
	\]

	\subsection{Solución}

	Para que una secuencia real tenga una DFT imaginaria pura, debe ser antisimétrica:

	\[
	x[n]
	=
	-x[N-n]
	\]

	con $N=8$.

	En el apartado a):

	\[
	x[0]=1
	\]

	por lo que no puede ser antisimétrica.

	En el apartado b):

	\[
	x[0]=1
	\]

	por lo que tampoco es antisimétrica.

	En el apartado c):

	\[
	x[0]=0
	\]

	y además:

	\[
	x[1]=-x[7]
	\]

	\[
	x[2]=-x[6]
	\]

	\[
	x[3]=-x[5]
	\]

	\[
	x[4]=0
	\]

	Por tanto, esta secuencia es antisimétrica.

	En el apartado d):

	\[
	x[1]=x[7],
	\qquad
	x[2]=x[6]
	\]

	por lo que es simétrica, no antisimétrica.

	Por tanto:

	\[
	\boxed{
	\text{La única secuencia cuya DFT es imaginaria pura es la c)}
	}
	\]

\section{Ejercicio 25}
	Utilizar la DFT para calcular el producto de los polinomios

	\[
	(1+x)
	\]

	y:

	\[
	(1+x+x^2)
	\]

	\subsection{Solución}

	Representamos los polinomios mediante sus coeficientes y completamos hasta longitud $N=4$:

	\[
	x[n]=[1,1,0,0]
	\]

	\[
	y[n]=[1,1,1,0]
	\]

	Calculamos sus DFT:

	\[
	\boxed{
	X[k]=[2,1-j,0,1+j]
	}
	\]

	\[
	\boxed{
	Y[k]=[3,-j,1,j]
	}
	\]

	Multiplicamos componente a componente:

	\[
	Z[k]
	=
	X[k]Y[k]
	\]

	\[
	Z[k]
	=
	[6,-1-j,0,-1+j]
	\]

	Aplicando la DFT inversa:

	\[
	z[n]
	=
	[1,2,2,1]
	\]

	Estos valores son los coeficientes del polinomio producto.

	Por tanto:

	\[
	\boxed{
	(1+x)(1+x+x^2)
	=
	1+2x+2x^2+x^3
	}
	\]

\section{Ejercicio 26}
	Calcular el producto de los números $321$ y $12$ utilizando la DFT.

	\subsection{Solución}

	Interpretamos los números como polinomios evaluados en $x=10$:

	\[
	321
	=
	3x^2+2x+1
	\]

	\[
	12
	=
	x+2
	\]

	Por tanto, las secuencias de coeficientes son:

	\[
	x[n]=[1,2,3,0]
	\]

	\[
	y[n]=[2,1,0,0]
	\]

	Calculamos sus DFT:

	\[
	\boxed{
	X[k]=[6,-2-2j,2,-2+2j]
	}
	\]

	\[
	\boxed{
	Y[k]=[3,2-j,1,2+j]
	}
	\]

	Multiplicamos componente a componente:

	\[
	Z[k]
	=
	X[k]Y[k]
	\]

	\[
	\boxed{
	Z[k]=[18,-6-2j,2,-6+2j]
	}
	\]

	Aplicando la DFT inversa:

	\[
	z[n]
	=
	[2,5,8,3]
	\]

	Por tanto, el polinomio producto es:

	\[
	3x^3+8x^2+5x+2
	\]

	Evaluando en $x=10$:

	\[
	3(10)^3+8(10)^2+5(10)+2
	\]

	\[
	=
	3000+800+50+2
	\]

	\[
	\boxed{
	321\cdot12=3852
	}
	\]

    
\end{document}